\section{\label{apx:calib}Hyperparameter calibration}
All tuning was performed on a held-out seed,
$99$,
disjoint from the grid seeds $0$--$2$;
on the clean condition only;
and applied unchanged to every condition.
Tuning per condition would have introduced exactly the confound the shared recipe exists to prevent.
The runs below are calibration measurements.

\paragraph{Peak learning rate.}
A lot size of $256$ makes $3$ epochs $192$ optimizer steps instead of $3072$,
so the rate inherited from a $B{=}16$ configuration is far too small:
at $2\times10^{-4}$ the clean condition reaches validation loss $0.9648$ against $0.8250$ at $B{=}16$,
and $6\times10^{-4}$ recovers only to $0.8700$.
Table~\ref{tab:lr} scans the peak rate on all three models.
We adopt $3.2\times10^{-3}$ for all three rather than one value per model,
because the data do not separate them.

\begin{table}[h]
\centering
\begin{tabular}{@{}lccc@{}}
\toprule
Peak LR & Qwen3-1.7B & Llama-3.2-1B & Falcon3-1B \\
\midrule
$1.0\times10^{-3}$ & 0.8425 & 1.0095 & 0.9402 \\
$2.0\times10^{-3}$ & 0.8186 & 0.9868 & 0.8931 \\
$\mathbf{3.2\times10^{-3}}$ & \textbf{0.8122} & \textbf{0.9809} & \textbf{0.8724} \\
$5.0\times10^{-3}$ & 0.8141 & 0.9876 & -- \\
$8.0\times10^{-3}$ & 0.8484 & -- & -- \\
\bottomrule
\end{tabular}
\caption{Final validation loss on the clean condition,
lot size $256$,
$3$ epochs,
calibration seed $99$.
The optimum is bracketed on both sides for Qwen and Llama;
the Falcon cell at $8\times10^{-3}$ was stopped once it trailed its $3.2\times10^{-3}$ counterpart at epoch~1.}
\label{tab:lr}
\end{table}

\paragraph{Where the search stops.}
Three seeds of the \textit{same} recipe differ at epoch $0$ by $0.0067$ ($0.8865$, $0.8886$, $0.8932$).
The gap between $3.2\times10^{-3}$ and $5\times10^{-3}$ in Table~\ref{tab:lr} is $0.0019$,
a quarter of that spread,
so the two settings are indistinguishable and any further point in this range would measure seed noise.
The same floor bounds what the design needs to resolve elsewhere:
the gap in validation loss between the non-private and private conditions at fixed learning rate is about $0.315$,
roughly $47\times$ the seed spread,
so the quantity the regression must estimate sits far above the noise level at three seeds per cell.

\paragraph{The tuned rate transfers to the private mechanism.}
Tuning on the clean condition carried the risk that a $16\times$ increase in learning rate would diverge against gradients clipped to nearly direction-only.
It does the opposite.
At $\varepsilon = 8$ the loss improves from $1.2808$ at $2\times10^{-4}$ to $1.1335$ at the tuned rate;
at $\varepsilon = 1$,
from $1.3712$ to $1.2532$ after one epoch.
The perturbed conditions gain as well ($2.8927 \to 2.7212$ after one epoch at $p{=}0.30$),
with no divergence and no NaN in any cell,
so no per-mechanism learning rate is required.

\paragraph{Two settings with flat response.}
Both were live hypotheses that would otherwise be re-proposed.
\textit{Clipping bound.} At the tuned rate and $\varepsilon = 8$,
$C = 0.5$ gives $1.1857$,
$C = 1.0$ gives $1.1859$ and $C = 2.0$ gives $1.2045$;
the first two agree to four decimals.
The knob is demonstrably live -- the fraction of clipped examples moves $0.684 / 0.645 / 0.432$ across the three -- it does not matter at this scale,
so $C = 1.0$ stands.
That fraction also falls as the rate rises ($0.773$ at $2\times10^{-4}$ against $0.645$ at $3.2\times10^{-3}$, both at $\varepsilon = 8$),
since larger steps shrink the gradients;
the two interact,
which is why $C$ was re-probed at the new rate rather than carried over.
\textbf{Decay onset.} Starting the cosine at $1.14$ epochs instead of $1.5$ changes Llama's final validation loss from $0.9809$ to $0.9775$ and Qwen's epoch-$0$ loss from $0.8865$ to $0.8856$.
Both sit below the resolution floor.

\paragraph{Sequence length.}
Truncation at $2048$ tokens is inexpensive and fixes real data loss:
at $1024$,
$50$ rows have prompts that fill the window entirely and $103$ completions are cut short;
at $2048$ those counts are $2$ and $0$. Doubling again to $4096$ drops nothing further while doubling activation memory,
which is the binding constraint in the private branch.

\paragraph{Validation loss serves as a diagnostic.}
An optional fixed-seed holdout logs token-weighted validation loss once per epoch,
splitting on a fixed generator rather than on the run seed so that every condition and seed holds out the same rows.
It is disabled for the main grid.
Each perturbed condition validates against its own perturbed distribution,
so the levels are incomparable by construction -- $0.96$,
$3.15$ and $4.81$ for the clean condition,
random replacement at $p{=}0.30$ and sanitization at $\varepsilon{=}3$ under otherwise identical settings.
That reflects a different target.
Within one fixed dataset it remains a valid training diagnostic,
which is how it is used above.
With the holdout disabled,
the grid trains on all $16{,}384$ rows and the privacy guarantee covers the whole dataset.

